\documentclass[11pt]{article}
\usepackage{amsmath,amssymb,amsthm}
\usepackage[margin=1in]{geometry}
\newtheorem{lemma}{Lemma}
\newtheorem{theorem}{Theorem}
\newtheorem{corollary}{Corollary}
\usepackage[hidelinks]{hyperref}
\newcommand{\Q}{\mathbb Q}
\begin{document}
\setlength{\emergencystretch}{1em}
\title{Three-adic balance and parity for ordinary quintic moments}
\author{}\date{}\maketitle
For the ordinary quintic support $\{2,5\}$, an exact mixed-adjoint
identity forces every root common to two consecutive moments to have
one specified three-adic valuation. When that valuation is $-1$, its
lower Newton cluster has width one. Real-root parity then excludes a
rational common factor. The only terminal progression left by this
argument is $N\equiv11\pmod{27}$. The \href{QUINTIC_SCALAR_MINOR_SFC.tex}{scalar-minor theorem} closes it.
Internal replay is separate from independent replay, formal verification,
external review and priority.

For $n\ge1$ put
\[
 P_n(t)=\sum_{k=0}^n\binom nk(2n+1)_{3k}t^k,\qquad
 D=5n+4,\qquad r=v_3(D).
\]
Every valuation is normalized by $v_3(3)=1$.
We work in a complete algebraically closed extension $K$ of $\Q_3$,
with valuation ring $\mathcal O$, maximal ideal $\mathfrak m$, and residue
field $k$. An overbar always means reduction modulo $\mathfrak m$;
it does not assert congruence modulo $3\mathcal O$.

\section{The exact mixed-adjoint input}
For every integer $m\ge0$ and parameter $\gamma$, define
\[
 T_{m,\gamma}(t)=\sum_{k=0}^m\binom mk(\gamma)_{3k}t^k,\qquad
 B_{0,\gamma}=1,\quad
 B_{j+1,\gamma}=(z+3j+\gamma+2)B_{j,\gamma}
                       -zB_{j,\gamma}' .
\]
Thus $P_n=T_{n,2n+1}$. Set
\[
 M=\gamma+3m,\quad A=(M+3)(M+4),\quad
 \kappa=\frac{3(m+1)}{A},
\]
\[
 B=B_{m,\gamma},\qquad C=B_{m+1,\gamma},\qquad
 D_*=B_{m+2,\gamma},\qquad H=D_*+\kappa z^3B.
\]
For $u\ne0$, write the cyclic column of a polynomial $J$ as
\[
 J(z)=J_0(z^3)+zJ_1(z^3)+z^2J_2(z^3),\qquad
 {\bf j}(u)=(J_0(u),J_1(u),J_2(u))^{\mathsf T}.
\]
\begin{lemma}[Exact imported mixed-adjoint equivalence]\label{lem:mixed-input}
For $m\ge1$, $\gamma>0$ and $u\ne0$, the two equalities
\[
 T_{m,\gamma}(u^{-1})=T_{m+1,\gamma+2}(u^{-1})=0
\]
are equivalent to $\operatorname{rank}[{\bf c}(u),{\bf h}(u)]<2$.
The equivalence is algebraic in characteristic zero, so it also holds
over $K$ for positive rational $\gamma$.
\end{lemma}
\begin{proof}[Precise imported identity and normalization]
This is Theorem \texttt{thm:gap-three-offset-}\allowbreak\texttt{two-rank} in the
\href{lower_exponent_initial_and_boundary.tex}{lower-exponent note}.

The proof of
\texttt{cor:gap-three-lower-}\allowbreak\texttt{two-eight-ninths}
there explicitly applies this algebraic identity over $K$.
The following formulas specify the imported normalization completely.

For $R=z^{3m}T_{m,\gamma}(z^{-3})$ and $E=M-zD_z$,
coefficient comparison gives
\[
 (\gamma)_2z^{3(m+1)}T_{m+1,\gamma+2}(z^{-3})
       =[z^3(E)_2+(E)_5]R.
\]
The fourth-order equation is
\[
 R^{(4)}-\frac{3(M-1)}zR'''
 +\frac{3M(M-1)}{z^2}R''
 -\left(\frac{M(M-1)(M+1)}{z^3}+1\right)R'
 +\frac{3m}zR=0.
\]
Reduction of the preceding operator yields $qR+bR'+cR''+dR'''$, where
\[
\begin{aligned}
 q&=(M)_5+[M(M+1)-6m(M+2)]z^3,\\
 b&=3z[-M(M+1)A+(m+1)z^3],\\
 c&=3Mz^2A,\qquad d=-z^3A.
\end{aligned}
\]
For $\xi^3=1$, the zero adjoint concomitants for
$f_\xi=e^{-\xi z}z^{1-M}B(\xi z)$, after imposing $R=0$
and eliminating $R'''$ with $d\ne0$, have coefficients
\[
 U_\xi=d(a_3f_\xi-f_\xi')-cf_\xi,\qquad
 V_\xi=d[f_\xi''-(a_3f_\xi)'+a_2f_\xi]-bf_\xi,
\]
where $a_3=-3(M-1)/z$, $a_2=3M(M-1)/z^2$.
If $g_\xi=e^{-\xi z}z^{1-M}$, the exact simplifications are
\[
 U_\xi/g_\xi=-z^2A C(\xi z),\qquad
 \bigl(V_\xi+(3M+5)U_\xi/z\bigr)/g_\xi=-zA H(\xi z).
\]
The full concomitant matrix has rank three. Its nullspace is spanned
by the nonzero jet $(R,R',R'',R''')$. This is the imported rank fact:
it makes elimination an equivalence, including multiple roots,
rather than just a necessary relation.
Fourier evaluation at $z,\omega z,\omega^2z$ is invertible when
$z^3=u\ne0$, and identifies the evaluation and cyclic-column ranks.
Every gauge has been removed from the resulting polynomial rank identity.
The source proof justifies this same algebraic identity over $K$.
No uniform rank assertion is imported.
\end{proof}

\section{Balance, the one-root cluster and parity}
\begin{lemma}[Every common root is balanced]\label{lem:balance}
If $P_n$ and $P_{n+1}$ have a common root $t\in K$, then
\[
 v_3(t)=1-r,\qquad
 \overline{t/3^{\,1-r}}=-\bigl(5\overline{D/3^r}\bigr)^{-1}\in k^\times.
\]
In particular a common root is impossible when $r\le1$.
\end{lemma}
\begin{proof}
Put $\gamma=2n+1$, $u=t^{-1}$ and choose $z^3=u$.
The constant terms are one, so $t,z,u$ are nonzero.
Every $(\gamma)_{3k}$ contains at least $k$ integer factors divisible
by three. Thus every coefficient of
\[
 z^{3n}P_n(z^{-3})
\]
at $z^{3(n-k)}$ has valuation at least $k$.
If $v_3(z)<1/3$, its leading term is uniquely of smallest valuation,
which is impossible at a root. Consequently $v_3(u)\ge1$.

Use Lemma~\ref{lem:mixed-input} with $m=n$:
\[
 B=B_{n,\gamma,3},\quad C=B_{n+1,\gamma,3},\quad
 D_*=B_{n+2,\gamma,3},\quad
 H=D_*+\kappa z^3B,\qquad
 \kappa=\frac3{5(5n+4)}.
\]
A common root forces the two cyclic columns $C,H$, evaluated at $u$,
to have rank less than two. Their coefficient vectors, as well as
that of $B$, are integral because $\gamma$ is an integer and
$v_3(u)>0$.
On reduction modulo the maximal ideal, $u$ is zero and the auxiliary
operator is
\[
 \mathsf A=z-zD_z+\bar\gamma-1
 \quad\hbox{on }k[z]/(z^3),\qquad \mathsf A^3=\mathsf A.
\]
The vectors $1,\mathsf A1,\mathsf A^21$ form a basis since they are
monic of degrees zero, one and two. Moreover $\bar\gamma\in\mathbb F_3$.
For every $n\ge1$, $\mathsf A^n1$ and $\mathsf A^{n+1}1$
are independent: $\mathsf A$ has the three distinct eigenvalues
$0,1,-1$, the cyclic vector has a nonzero component in every
eigenspace, and the two nonzero eigenspaces give a Vandermonde minor.
The reduced columns $B,C,D_*$ are respectively
$\mathsf A^n1,\mathsf A^{n+1}1,\mathsf A^{n+2}1$.

Write $\eta=\kappa u$. If $v_3(\eta)>0$, the pair $C,H$ reduces
to $\mathsf A^{n+1}1,\mathsf A^{n+2}1$, of rank two.
If $v_3(\eta)<0$, first multiply the second column by $\eta^{-1}$;
it then reduces to the independent pair
$\mathsf A^{n+1}1,\mathsf A^n1$.
Both contradict rank deficiency. Hence $v_3(\eta)=0$.

Since $n\ge1$, $\mathsf A^{n+2}1=\mathsf A^n1$.
The pair now reduces to
$\mathsf A^{n+1}1,(1+\bar\eta)\mathsf A^n1$.
It has rank less than two exactly when $\bar\eta=-1$.
As $v_3(\kappa)=1-r$, balance gives $v_3(u)=r-1$
and $v_3(t)=1-r$.
Writing $t=3^{1-r}w$, one has
$\eta=1/[5(D/3^r)w]$, proving the residue assertion.
The inequality $v_3(u)\ge1$ excludes $r\le1$.
\end{proof}

\begin{lemma}[A unique valuation-minus-one root]\label{lem:one-root}
If $n\equiv1\pmod9$, the lower Newton polygon of $P_n$ has first
edge $(0,0)$ to $(1,1)$, of width one, and all other slopes
strictly exceed one. Thus it has exactly one root of valuation $-1$.
\end{lemma}
\begin{proof}
Write $2n+1=3g$. Then $g\equiv1\pmod3$ and, exactly,
\[
 v_3\!\left(\binom nk(2n+1)_{3k}\right)
 =k+v_3\!\left(\binom nk(g)_k\right).
\]
Indeed the integer multiples of three in the rising factorial
are precisely $3g,3(g+1),\ldots,3(g+k-1)$.
For $k=1$, both $n$ and $g$ are units, giving valuation one.
For $k=2$, $\binom n2=n(n-1)/2$ is divisible by nine.
For $k\ge3$, $(g)_k$ contains a multiple of three.
All coefficients at $k\ge2$ therefore have valuation strictly
larger than $k$. Together with $c_0=1$, this proves the stated
Newton polygon and its width-one root count.
\end{proof}

\begin{theorem}\label{thm:new-ray}
For every integer $n\ge1$ with $n\not\equiv10\pmod{27}$,
\[
 \gcd(P_n,P_{n+1})=1
\]
over $\Q$ and over $\mathbb C$.
Equivalently, the only remaining ordinary terminal orders for
support $\{2,5\}$ lie on $N\equiv11\pmod{27}$.
\end{theorem}
\begin{proof}
If $v_3(5n+4)\le1$, Lemma~\ref{lem:balance} excludes any common
root over $K$, hence the rational resultant is nonzero.
Suppose instead $r=2$; this implies $n\equiv1\pmod9$.
Every common root must have valuation $-1$ by
Lemma~\ref{lem:balance}, and the lower polynomial has exactly one
such root, with multiplicity one, by Lemma~\ref{lem:one-root}.
The degree of the rational gcd is therefore at most one.

On the other hand it cannot have degree one. The even-index member
is positive at every real argument, by its Gamma-moment integral.
Every common real-coefficient factor has no real roots and therefore
has even degree. Indeed for real $t$ the finite Gamma-moment expansion is
\[
 P_j(t)=\frac1{\Gamma(2j+1)}
       \int_0^\infty x^{2j}(1+tx^3)^j e^{-x}\,dx.
\]
When $j$ is even its integrand is nonnegative and positive on an
interval, so $P_j(t)>0$. This supplies the parity statement directly.
A nonconstant rational gcd would have even degree at least two,
a contradiction.

The only possibility not excluded is $r\ge3$, namely
$5n+4\equiv0\pmod{27}$ or $n\equiv10\pmod{27}$.
A rational Bezout identity remains valid over $\mathbb C$.
For $AB\ne0$ the exact normalization is
\[
 \mathcal L_1\bigl((X^2(A+BX^3))^j\bigr)
       =A^j(2j)!P_j(B/A).
\]
If $A=0$ or $B=0$, every positive-order moment is nonzero because
the phase is a nonzero monomial.
\end{proof}

\begin{corollary}[The remaining degree-five target]\label{cor:degree-five-ray}
Strong Factorial for every ordinary univariate binomial of degree
at most five is equivalent to
\[
 \gcd(P_{N-1},P_N)=1
 \quad\hbox{for every }N\ge11,\qquad N\equiv11\pmod{27}.
\]
\end{corollary}
\begin{proof}
Corollary \texttt{cor:ordinary-quintic-}\allowbreak\texttt{ray-reduction} in the
\href{lower_exponent_initial_and_boundary.tex}{lower-exponent note}
reduces the degree-five problem to the same support at
$N\ge11$, $N\equiv2\pmod9$.
Theorem~\ref{thm:new-ray} removes the two residue classes
$N\equiv2,20\pmod{27}$.
The remaining tests are part of the original degree-five assertion,
which proves the converse. No claim discharges them.
\end{proof}

\begin{corollary}[Two thirds of the sparse no-carry-failure family]\label{cor:sparse-family}
For $n=2^{12s+6}$, the new theorem excludes the pair whenever
$s\not\equiv0\pmod3$.
\end{corollary}
\begin{proof}
The powers of two modulo $27$ have period $18$.
As $s$ runs modulo three, $n$ is respectively $10,1,19$ modulo $27$.
Only the first is retained by Theorem~\ref{thm:new-ray}.
\end{proof}
\begin{corollary}[The interpolated common index is not an ordinary order]\label{cor:interpolated-index}
Let $(\tau_*,w_*)$ be the unique integral common point in
Theorem \texttt{thm:p3-index-}\allowbreak\texttt{interpolation-}\allowbreak\texttt{obstruction} of the
\href{lower_exponent_initial_and_boundary.tex}{lower-exponent note}.
Then neither $\tau_*$ nor $m_*=1+27\tau_*$ is a nonnegative ordinary
integer. This asserts no irrationality, transcendence or negative-integer
exclusion.
\end{corollary}
\begin{proof}
That theorem gives $\tau_*\in\mathbb Z_3$ and a unit $w_*\in\mathbb Z_3$,
with the two interpolated series zero. If $\tau_*$ were a nonnegative
integer, $m_*=1+27\tau_*$ would be a positive integer and the series
would terminate, exactly as
\[
 P_{m_*}(w_*/3)=P_{m_*+1}(w_*/3)=0.
\]
But $m_*\equiv1\pmod{27}$ contradicts Theorem~\ref{thm:new-ray}.
If instead $m_*$ were a nonnegative integer, it is at least one
because $m_*\equiv1\pmod{27}$, and
$\tau_*=(m_*-1)/27$ would be a nonnegative integer, giving the same
contradiction. The finite-degree and parity arguments used here
are not assertions about noninteger interpolation indices.
\end{proof}

\section{Scope of the replay}
The exact checker verifies the displayed Ore reduction symbolically
in $m,M$ and the two gauge-stripped coefficient identities with a
formal polynomial jet. It checks all three reductions of $\gamma$
and both positive-power parity classes in the cyclic operator,
including every residue of $\eta$.
Literal coefficient valuations, a few direct gcd controls and
the full exponent-period calculation support the formulas.
The all-index conclusion follows from the proofs, not a finite
index search. \href{QUINTIC_SCALAR_MINOR_SFC.tex}{the scalar-minor theorem} closes the $N\equiv11\pmod{27}$ progression
and all ordinary binomials through degree five. Unrestricted ordinary
binomial SFC and unrestricted FC/SFC remain open.
\end{document}
